\section{Introduction}
\label{sec:intro}

When a particle-in-cell (PIC) simulation starts, the bunch of particles must
carry an initial electromagnetic field, or the first push already acts on the
wrong force. The field of a moving bunch in free space is known in closed form
for each point charge, but a simulation reads that field on a staggered mesh
with thousands of observation points, so it is computed by a fast sum over the
source charges rather than by evaluating the closed form everywhere. We do that
sum as a TSC deposit onto a zero-padded grid, a sampled exact TSC-convolved
kernel, and a free-space FFT convolution, with staggered Yee read-out. This
paper gives the exact discrete recipe a re-implementer needs, and it measures
the cost and the accuracy of that recipe on a single V100, naming the reference
for every error statement.

The established workaround avoids computing the initial field altogether: one
starts from a field-free state and marches the bunch forward, letting the
fields build up in the first few steps. The present-time field of a uniformly
moving bunch has no transient, so this history has no physical basis. We
measure how well the discrete method reproduces that field and how much it
costs, and we state the conditions under which it fails
(Section~\ref{sec:limits}).
